%==================== CHAPTER 3 ====================
\ChapterPage{Proposed System}

\section{Problem Statement}
Today's tools were designed for an era when information was stored in structured databases and workflows were governed by hard-coded rules. Posting a job meant setting up a job and filling out a spreadsheet of applications, then individually reading through them. Marketing involved developing a template and filling the same in for every channel. Research involved having a number of browser windows and making notes. As the amount of work has increased significantly, and the risks of failure have become more obvious, these processes have not significantly changed.

A recruiter using a keyword filter who looks at 1,000 resumes will overlook good candidates that did not state their experience in the proper manner. A screening system without scores to explain them means that the recruiter will not be able to detect errors or bias. Without a review process, a marketing team can find itself on the verge of a brand crisis after publishing just one bad post. If a podcast host puts together a guest briefing without verifying information from unreliable sources, then he could be spreading false information or defamation on the air. In each instance, the information the team is after is stored in a different type of application or system than the one they're searching, each with its own login, and each with its own search methods.
What is missing is a system that takes over the repetitive, high-volume parts of these workflows, explains every output clearly enough that a person can check it, and then waits for that person's approval before doing anything that cannot be undone.

\section{Scope}
This project covers four business workstreams, defined in collaboration with our industry partner \Partner: HR recruitment and HR assistance, marketing content operations, podcast guest background research, and enterprise integrations with Zoho, Workday and Greenhouse-style platforms. In all four, a common workflow router, a shared retrieval-augmented knowledge layer, approval gates at every “consequential” step to be filled by humans, and a complete audit log, with a privacy protection of PII, all of this as a reusable framework.

Clear demarcation of what the system does and doesn't do. It helps people, it doesn't make hiring decisions, rejection decisions, salary decisions or disciplinary decisions on their behalf. It stores content that is for approval and it publishes nothing that hasn't been approved. All development is done using synthetic data, public data or sanitised data, and not live production credentials, or sandbox APIs. It is planned to be expanded to other business functions later, for example finance and sales; however, this is not covered by the current scope.
\section{Proposed System}
\textbf{\ProjectTitle} is a modular multi-agent platform built around a single idea: business workflows should be fast and consistent where the work is repetitive, and human where the decisions matter.

\vspace{12pt}
\RunIn{System Overview}
All interactions begin at a Web dashboard. A request is triggered by a user or by an incoming business events and a workflow router categorizes it and routes it to the appropriate set of agents. Those agents retrieve information from a shared context store and invoke enterprise connectors as necessary. Any action with a consequence -- sending an email to a candidate, publishing a post, writing a record back to an HR system -- is put into an approval queue, and is only executed once it is approved. It is examined by a human being, approved and/or edited by a human, and only then does the system act. All exchanges are logged in an audit log.
\vspace{12pt}
\RunIn{Key Components}
\begin{itemize}
  \item \textbf{Web Dashboard} -- A Next.js interface where HR managers, marketing teams and researchers interact with the system: creating jobs, reviewing candidate rankings, editing content drafts and clearing the approval queue. Candidates have a separate public-facing apply page.
  \item \textbf{Workflow Router} -- Classifies each incoming request as an HR, marketing, podcast research or integration task and hands it to the matching set of agents.
  \item \textbf{HR Agents} -- A set of specialised agents covering rubric generation from job descriptions, resume parsing and adversarial scanning, claim verification, evidence-backed screening, interview question generation, live AI interviewing, candidate ranking and follow-up. Separate assistants handle employee queries about onboarding, leave, attendance and HR policy.
  \item \textbf{Marketing Agents} -- Agents that turn a theme or source content into a content calendar, campaign and creative briefs, audience segments and channel-specific social copy. All drafts wait in an approval queue before going anywhere.
  \item \textbf{Podcast Research Agents} -- Agents that gather public sources about a guest, index them for retrieval, build a structured briefing and label every claim as a verified fact, an allegation or an opinion. Sensitive-topic guardrails are applied before the briefing reaches the host.
  \item \textbf{Integration Agent and Connector Layer} -- A unified connector interface for Zoho People, Zoho CRM and Recruit, Workday-style HRIS, Greenhouse-style ATS, CRM and email, calendar and social media APIs. Handles OAuth, rate limits, pagination, webhooks, retries and idempotency so that agents do not need to know the details of each system.
  \item \textbf{Knowledge (RAG) Layer} -- Documentation of the company policies, brand guidelines, job descriptions, approved rubrics, previous campaigns, guest research sources and more in a vector store, agents always look up in this store, rather than relying on just 'memory'.
  \item \textbf{Compliance Reviewer and Guardrails} -- Acts as the bridge between the output from each agent and the approval queue. Looks for PII that shouldn't be part of a prompt or log, for prompt injection in user-contributed content, and for unfounded claims with a cited source. Raises these issues for human correction rather than letting them go unnoticed.
  \item \textbf{Approval Layer} -- It holds the rubrics, hiring decisions and drafts of contents until a human acts on them. Nothing goes ahead of this step without approval.
  \item \textbf{Audit Layer and Data Store} -- Supabase storage for every job, candidate, score, draft and report, including a versioned log of every prompt, tool call, source, draft, reviewer and final output. Secured by role-based access controls.
\end{itemize}

\vspace{12pt}
\RunIn{Workflow of the Proposed System}
Upon receipt of the request, the workflow router categorize it and route it to the right specialists. The specialists get the required information from the knowledge base and call the enterprise connector where required. All outputs that have an implication or impact are not immediately picked up but stored in the approval queue. The compliance reviewer examines it for PII leaking, prompt injection and unsupported claims before it is passed to the reviewer. The reviewer human analyzes it, approves, edits or declines it. Approved actions are sent through the connector layer. The audit log records every step from initial prompt to final output.